\paragraph{What the measurement model certifies.}
Table~\ref{tab:measurement-assumptions} defines the stored-record and
bounded-error models used below.

\begin{table}[htbp]
\centering\small
\caption{Targets and observation assumptions in the application. Error
allowances and hourly caps are sensitivity inputs, not verified device limits.}
\label{tab:measurement-assumptions}
\begin{tabular}{>{\raggedright\arraybackslash}p{0.17\linewidth}>{\raggedright\arraybackslash}p{0.39\linewidth}>{\raggedright\arraybackslash}p{0.34\linewidth}}
\toprule
Model & Target and observations & Scope of the guarantee\\
\midrule
Exact stored samples
& Targets are increments of stored cumulative samples. Retained endpoints
and queried internal samples are exact in this record model.
& Loss-contrast intervals on the stored sample grid.\\
Bounded internal errors
& The latent record includes one fixed internal-error vector with
$|e_k|\leq\delta E$, where $E$ is the retained block total. Endpoints, and hence the
retained total, remain exact. Errors are fixed before policy randomization.
& All records and fixed error vectors within these bounds.\\
Physical energy
& Actual energy and meter readings would require defensible bounds for
internal measurements, endpoints and any physical capacity constraints.
& We provide no verified physical calibration. Source meter-class metadata
alone does not establish the pointwise error model or exact endpoints.\\
\bottomrule
\end{tabular}
\end{table}

\paragraph{Unchanged precision under capacity controls.}
We retain the uncapped exact-total initial width $W_{r,0}$ and the target
$W_r\leq W_{r,0}/2$, so every shared capacity control tests the same
absolute uncertainty budget. The scalar control in Online Resource~1, S6
also reports results relative to the capped initial widths.
